\chapter{Products of coherent boundary operations}
\label{ch:mc-products}

Two central operations can be composed. Their product must retain
the homotopies attached to boundary maps, as well as their values
on individual boundaries. Splitting a string of boundary maps
between the two operations gives the required multiplication.
The diagonal contraction extends to this multiplication and
therefore compares every product of bulk operators.

The higher terms can be nonzero at arbitrarily large arity.
In one coordinate they are repeated divided differences.
Their diagonal values are derivatives, and their generating
series give coordinate translation. Thus the homotopies that
compare products retain dependence on the polynomial coordinate.

\section{Multiplying cochains by splitting a string}
\label{sec:mc-convolution}

Let \(n\geq1\), let \(k\) be a field of characteristic zero,
let \(R=k[x_1,\ldots,x_n]\),
and fix \(W\in R\) and \(c\in k\). Let \(\mathcal C\) be the
category of finite free based matrix factorizations of \(W-c\)
from Section~\ref{sec:cc-coefficients}.
Write \(N(E)=E(x)\) for the original boundary and
\(M(E)=\Delta_W\otimes_{k[y]}E(y)\) for the diagonal action.
Both are \(k\)-linear differential graded functors into
factorizations over \(R\). The modules \(M(E)\) may have
infinite rank. Their polynomial vectors are finite sums.

For \(X,Y\in\{N,M\}\), let \(C(X,Y)\) consist of families \(u\)
whose component on a composable string is a \(k\)-multilinear map
\[
 u_m:(s\mathcal C(E_1,E_0))\otimes_k\cdots\otimes_k
          (s\mathcal C(E_m,E_{m-1}))
       \longrightarrow\operatorname{Hom}_R(X(E_m),Y(E_0)).
\]
The suspension has parity \(|sf|=|f|+1\).
If \(u\) has parity \(p\), its value on the word
\(w=sf_1\otimes\cdots\otimes sf_m\) has parity \(p+|w|\),
where \(|w|=\sum_j(|f_j|+1)\).
The case \(m=0\) gives a map \(X(E)\to Y(E)\) at each object.
Families are a product over all arities and object strings.

For \(X,Y,Z\in\{N,M\}\), \(u\in C(Y,Z)\), and
\(v\in C(X,Y)\), define their
convolution product by
\begin{equation}
 (u\star v)_m(w)
 =\sum_{j=0}^m
 (-1)^{|v|\sum_{i=1}^j(|f_i|+1)}
 u_j(sf_1,\ldots,sf_j)\,
 v_{m-j}(sf_{j+1},\ldots,sf_m).
 \label{eq:mc-cup}
\end{equation}
The intermediate object is \(E_j\). Thus every composition
in the sum has source \(X(E_m)\) and target \(Z(E_0)\).
The sign moves the map \(v\) past the first \(j\) suspended
inputs. For a word of length zero, this is ordinary composition.
For one input, it is
\[
 (u\star v)_1(sf)
    =u_0v_1(sf)+(-1)^{|v|(|f|+1)}u_1(sf)v_0.
\]
This formula already shows why the product of the nullary
values does not determine the full product.

Convolution is associative. For three composable cochains
\(u,v,z\), both parenthesizations sum over the same three
consecutive blocks. If the first two blocks have parities
\(a,b\), the sign in either order is
\[
 (-1)^{|v|a+|z|(a+b)}
\]
This proves associativity term by term. The identity has the
identity map in arity zero and zero in every positive arity.

Retain the word coderivation \(b\) of
Section~\ref{sec:cc-cochains}:
\[
 b_1(sf)=s\,\delta f,\qquad
 b_2(sf,sg)=(-1)^{|f|}s(fg),\qquad b^2=0.
\]
For \(u\in C(X,Y)\) of parity \(p\), put
\[
 (d_0u)(w)=\delta(u(w))+(-1)^p u(bw).
\]
The first term is the map differential
\[
 \delta(u(w))
    =D_{Y(E_0)}u(w)-(-1)^{p+|w|}u(w)D_{X(E_m)}.
\]
Its square is zero because both factorizations have curvature
\(W(x)-c\). The input term uses the differential \(-b\) on
words. Consequently \(d_0^2=0\): the two mixed terms cancel,
and the remaining terms are \(\delta^2u-u b^2\).
The coderivation rule for \(b\) and the product rule for map
differentials give
\[
 d_0(u\star v)=d_0u\star v+(-1)^{|u|}u\star d_0v.
\]
Each term differentiates one of the two blocks in
\eqref{eq:mc-cup}, with the displayed tensor sign.

\section{The differential of a coherent transformation}
\label{sec:mc-differential}

For a functor \(X\in\{N,M\}\), define an odd cochain
\(\tau_X\in C(X,X)\) by
\[
 (\tau_X)_1(sf)=X(f),\qquad(\tau_X)_m=0\quad(m\ne1).
\]
The functor respects differentials and composition.
The arity-one component of \(d_0\tau_X\) is zero.
In arity two,
\[
 (d_0\tau_X)_2(sf,sg)
   =-(-1)^{|f|}X(fg)
   =(\tau_X\star\tau_X)_2(sf,sg).
\]
All other arities vanish. Thus
\begin{equation}
 d_0\tau_X=\tau_X\star\tau_X.
 \label{eq:mc-twisting}
\end{equation}
Define the differential on \(C(X,Y)\) by
\begin{equation}
 du=d_0u-\tau_Y\star u+(-1)^{|u|}u\star\tau_X.
 \label{eq:mc-total}
\end{equation}
Expansion gives
\[
 d^2u=-(d_0\tau_Y-\tau_Y^2)\star u
                   +u\star(d_0\tau_X-\tau_X^2)=0,
\]
where \(\tau_X^2=\tau_X\star\tau_X\).
The same expansion gives the product rule for \(d\).
The spaces \(C(X,Y)\), with this differential and convolution,
therefore form a differential graded category.

The endomorphism algebra
\[
 \mathcal Z=C(N,N)
\]
is the Hochschild cup algebra used here. Its elements have
unsuspended outputs. If \(U_m=s\,u_m\), the cochain \(U\) has
parity \(|u|+1\), and
\[
 \mathfrak dU=s\,du
\]
for the differential of Section~\ref{sec:cc-cochains}.
Indeed, the input terms in both formulas are
\((-1)^{|u|}u b\). The two convolution terms in
\eqref{eq:mc-total} are exactly the two outer compositions in
the coderivation commutator. This identifies the cochain
complexes and fixes the multiplication on coherent operations.

\section{Lifting the diagonal contraction}
\label{sec:mc-lift}

The maps \(\pi_E,i_E,h_E\) of Theorem~\ref{thm:da-unit}
define cochains \(\pi,i,h\) supported in arity zero.
Their sources and targets are
\[
 \pi\in C(M,N),\qquad i\in C(N,M),\qquad h\in C(M,M).
\]
Their parities are \(0,0,1\), respectively. For \(d_0\) they
satisfy
\[
 d_0\pi=d_0i=0,\quad d_0h=1-i\pi,\quad
 \pi i=1,\quad \pi h=hi=h^2=0.
\]
Products of these nullary cochains are ordinary compositions.
Naturality of evaluation is the further identity
\[
 \pi\star\tau_M=\tau_N\star\pi.
\]
The inclusion \(i\) is generally not closed for the full
differential \(d\). Its failure is the difference between
\(M(f)i\) and \(iN(f)\). Inserting \(h\) successively corrects
that difference.

Filter cochains by their least possible arity. The cochain
\(h\star\tau_M\) has arity one, so
\begin{equation}
 K=\sum_{j\geq0}(h\star\tau_M)^{\star j},\qquad
 \mathcal I=K\star i,\qquad \mathcal H=K\star h
 \label{eq:mc-lifted-maps}
\end{equation}
are well-defined: only finitely many terms contribute at each
arity. The maps \(\mathcal I,\mathcal H\) have the same sources,
targets and parities as \(i,h\).
No sum of vectors of infinite polynomial degree is taken.

\begin{theorem}[The contraction among coherent transformations]
\label{thm:mc-contraction}
For the full differential \eqref{eq:mc-total},
\[
 d\pi=d\mathcal I=0,\qquad
 d\mathcal H=1-\mathcal I\star\pi,\qquad
 \pi\star\mathcal I=1,
\]
and
\[
 \pi\star\mathcal H=0,\qquad
 \mathcal H\star\mathcal I=0,\qquad
 \mathcal H\star\mathcal H=0.
\]
\end{theorem}
\begin{proof}
Suppress the convolution symbol, and write \(t=\tau_M\),
\(s_0=\tau_N\). The arity filtration gives
\(K=(1-ht)^{-1}\). Since \(\pi h=h^2=0\),
\[
 \pi K=\pi,\qquad hK=h,\qquad K=1+\mathcal Ht.
\]
Differentiate \(K(1-ht)=1\). The parity of \(K\) is zero, so
\[
 d_0K=K\,d_0(ht)\,K.
\]
Using \eqref{eq:mc-twisting} and \(\pi t=s_0\pi\),
\[
 d_0(ht)=(1-i\pi)t-ht^2
          =(1-ht)t-i s_0\pi.
\]
It follows that
\[
 d_0K=tK-\mathcal I s_0\pi,\qquad
 d_0\mathcal I=t\mathcal I-\mathcal I s_0,
\]
and
\[
 d_0\mathcal H
   =t\mathcal H+K-\mathcal I\pi
   =t\mathcal H+\mathcal Ht+1-\mathcal I\pi.
\]
The identities \(\pi i=1\) and \(\pi h=0\) were used in these
two multiplications. Formula~\eqref{eq:mc-total} now gives
\(d\mathcal I=0\) and \(d\mathcal H=1-\mathcal I\pi\).
It gives \(d\pi=0\) from naturality.
Finally,
\[
 \begin{gathered}
 \pi\mathcal I=\pi Ki=\pi i=1,\qquad
 \pi\mathcal H=\pi Kh=0,\\
 \mathcal H\mathcal I=KhKi=Khi=0,\qquad
 \mathcal H^2=KhKh=Kh^2=0.
 \end{gathered}
\]
Every calculation is valid at each fixed arity, where its sums
are finite.
\end{proof}

These cochains have explicit components. For a string
\(f_j:E_j\to E_{j-1}\), write
\(\sigma=\sum_j(|f_j|+1)\). Then
\[
 \mathcal I_m(sf_1,\ldots,sf_m)
   =h_{E_0}M(f_1)h_{E_1}\cdots
                h_{E_{m-1}}M(f_m)i_{E_m}.
\]
The homotopy component is
\[
 \mathcal H_m(sf_1,\ldots,sf_m)
   =(-1)^\sigma h_{E_0}M(f_1)h_{E_1}\cdots
                h_{E_{m-1}}M(f_m)h_{E_m}.
\]
The sign in \(\mathcal H_m\) moves its final odd \(h\) past the
suspended input word. In arity zero these formulas give \(i,h\).

\section{The multiplicative comparison at every arity}
\label{sec:mc-all-arities}

Let \(B_W=\operatorname{End}_{k[x,y]}(\Delta_W)\), with its
map differential. For \(a\in B_W\), the operator
\(L(a)=a\otimes1_E\), at each boundary \(E\), is a nullary
cochain in \(C(M,M)\). It satisfies
\[
 L(ab)=L(a)\star L(b),\qquad dL(a)=L(\delta a),\qquad L(1)=1.
\]
For the differential equation, the only positive-arity term is
\(-\tau_M\star L(a)+(-1)^{|a|}L(a)\star\tau_M\).
It vanishes by the tensor sign rule
\(L(a)M(f)=(-1)^{|a||f|}M(f)L(a)\).
Thus \(L\) is a strict differential graded algebra action.

Define maps \(\Psi_r:(sB_W)^{\otimes_R r}\to s\mathcal Z\),
for \(r\geq1\), by
\begin{equation}
 \Psi_r(sa_1,\ldots,sa_r)
   =s\bigl(\pi\star L(a_1)\star\mathcal H\star L(a_2)
           \star\cdots\star\mathcal H\star L(a_r)
           \star\mathcal I\bigr).
 \label{eq:mc-comparison}
\end{equation}
There are \(r-1\) factors \(\mathcal H\). The unsuspended
cochain has parity \(\sum_j|a_j|+r-1\), so \(\Psi_r\) has
parity zero on suspended inputs.

\begin{theorem}[Products of diagonal operators]
\label{thm:mc-products}
The maps \eqref{eq:mc-comparison} form a strictly unital
\(R\)-linear \(A_\infty\)-morphism from \(B_W\) to the
Hochschild cup algebra \((\mathcal Z,d,\star)\).
Its first component is the coherent cochain constructed in
Theorem~\ref{thm:cc-coherent}.
At each bulk arity \(r\) and boundary arity \(m\), the formula
is a finite sum of polynomial operators with their specified
homotopies.
Every positive boundary-arity component vanishes on an
identity boundary input.
\end{theorem}
\begin{proof}
For a differential graded algebra \(A\), use the suspended
word differential
\[
 b_{A,1}(sa)=s\,d_Aa,\qquad
 b_{A,2}(sa,sb)=(-1)^{|a|}s(ab).
\]
The components \(\Psi_r\) extend to a coalgebra map by
partitioning each bulk word into consecutive nonempty blocks,
as in Section~\ref{sec:da-boundary-operators}.

Differentiate the unsuspended expression in
\eqref{eq:mc-comparison}. Differentiating \(L(a_j)\) has sign
\[
 (-1)^{\sum_{i<j}|a_i|+j-1}.
\]
These are the terms obtained by applying \(b_{B_W,1}\) to
the input word. Differentiating the homotopy after \(a_j\)
has sign
\[
 (-1)^{\sum_{i\leq j}|a_i|+j-1}
\]
and replaces \(\mathcal H\) by \(1-\mathcal I\star\pi\).
The term \(1\) merges the adjacent bulk operators.
The term \(-\mathcal I\star\pi\) splits the output into the
convolution product of the two corresponding shorter components.
The displayed sign is the source multiplication sign and
the negative of the target multiplication sign, respectively.
The closedness of \(\pi,\mathcal I\) gives no other terms.
This proves the length-one component of
\[
 b_{\mathcal Z}\Psi=\Psi b_{B_W}.
\]
The difference is a coderivation relative to \(\Psi\).
If it vanishes on shorter words, its value on the next word
has zero reduced coproduct and therefore has length one.
The component just proved is zero. Induction proves the full
coalgebra equation.

For an initial bulk unit in a longer word, the formula
contains \(\pi\mathcal H=0\). A final unit gives
\(\mathcal H\mathcal I=0\), and an interior unit gives
\(\mathcal H^2=0\). For \(r=1\), the unit gives
\(\pi\mathcal I=1\).
The maps are \(R\)-linear in the left bulk coefficients
because every contraction map is \(R\)-linear.

For \(r=1\), inserting the components of \(\mathcal I\)
gives exactly \eqref{eq:cc-components}.
For fixed boundary arity \(m\), convolution distributes the
boundary word among \(r-1\) factors \(\mathcal H\) and the
last factor \(\mathcal I\). Only finitely many consecutive
block decompositions occur. Each summand has \(m+r-1\)
ordinary homotopies \(h_E\), proving the finite-arity assertion.
An identity boundary input lies in a component of
\(\mathcal H\) or \(\mathcal I\). In \(\mathcal H\), it gives
two adjacent \(h\)'s. In \(\mathcal I\), it gives either
two adjacent \(h\)'s or a final \(hi\). These products vanish,
so each convolution summand vanishes on that input.
\end{proof}

For each boundary \(E\), evaluation at arity zero defines
\[
 \operatorname{ev}_E:\mathcal Z\longrightarrow
           \operatorname{End}_R(E(x)),\qquad
 \operatorname{ev}_E(u)=u_0(E).
\]
This is a differential graded algebra map: arity-zero
convolution is ordinary composition, and its differential
is the map differential. On the suspended components,
\[
 s\operatorname{ev}_Es^{-1}\Psi_r(sa_1,\ldots,sa_r)
   =s\bigl(\pi_E L_E(a_1)h_E\cdots
                     h_E L_E(a_r)i_E\bigr).
\]
Thus evaluating the coherent comparison recovers exactly the
boundary operator comparison of Theorem~\ref{thm:da-all-arity}.
The two maps in
\[
 B_W\xrightarrow{\ \Psi\ }\mathcal Z
       \xrightarrow{\ \operatorname{ev}_E\ }
                    \operatorname{End}_R(E(x))
\]
have different levels: the first is an \(A_\infty\)-morphism
and the second is a strict differential graded algebra map.
The cubic class of Theorem~\ref{thm:cc-cubic} is nonzero
before evaluation and zero on every resulting object
cohomology. Its compatibility data cannot be recovered
from those object cohomologies alone.

\section{Divided differences and coordinate translation}
\label{sec:mc-translation}

In one coordinate the homotopy has the form
\[
 h(v+\theta w)=\theta\,\frac{v(x,y)-v(x,x)}{x-y}.
\]
It vanishes on the exterior summand, and every \(M(f)\)
preserves that summand. Thus \(hM(f)h=0\).
Consequently \(\mathcal H\) has only its nullary component,
and \(\mathcal I\) has only nullary and unary components.
Every \(\Psi_r\) therefore has boundary arity at most one.

Let \(\iota\) be the odd diagonal operator sending
\(\theta\) to \(1\) and \(1\) to zero.
For a polynomial matrix \(F(x,y)\), put
\[
 T(F)=\frac{F(x,y)-F(x,x)}{x-y}.
\]
Write \(f^{(r)}\) for the \(r\)-th entrywise derivative.
For \(F=f(y)\), expansion of a monomial
\(y^d=(x+(y-x))^d\) gives
\[
 \left.T^r(f(y))\right|_{y=x}
       =\frac{(-1)^r}{r!}f^{(r)}(x).
\]
Indeed, each application removes the constant term in \(y-x\)
and divides the remainder by \(-(y-x)\). After \(r\)
applications, evaluation extracts the coefficient of
\((y-x)^r\), with sign \((-1)^r\). Linearity proves the
formula for every polynomial matrix.

Write \(\psi_r=s^{-1}\Psi_r(s\iota,\ldots,s\iota)\).
For every boundary \(E\) and boundary map \(f\),
\begin{equation}
 (\psi_r)_0(E)=\frac{(-1)^{r-1}}{r!}D_E^{(r)}(x),\qquad
 (\psi_r)_1(sf)=\frac{(-1)^r}{r!}f^{(r)}(x),\qquad
 (\psi_r)_m=0\quad(m\geq2).
 \label{eq:mc-derivatives}
\end{equation}
To prove the first formula, the rightmost \(\iota\) selects
the divided-difference matrix
\(A=(D_E(x)-D_E(y))/(x-y)=-T(D_E(y))\) in \(i_E\).
Each remaining pair \(h,\iota\) applies \(T\) again.
For the second formula, \(\mathcal I_1(sf)=hM(f)i\) starts
with \(T(f(y))\), and the \(r\) copies of \(\iota\) give
\(r\) applications of \(T\). The vanishing in higher
boundary arity follows from \(hM(f)h=0\).

For the cubic boundary
\[
 D_E(x)=\begin{pmatrix}0&x\\x^2&0\end{pmatrix},
\]
the binary bulk component has
\[
 (\psi_2)_0(E)=-\begin{pmatrix}0&0\\1&0\end{pmatrix},
 \qquad
 (\psi_2)_1(s(x^2\,1_E))=1_E.
\]
More generally,
\((\psi_r)_1(s(x^r1_E))=(-1)^r1_E\), so the mixed
components are nonzero for every bulk arity \(r\).
The operator \(\iota\) is not closed:
\(\delta\iota=(W(x)-W(y))/(x-y)\).
Its differential remains in the identities of
Theorem~\ref{thm:mc-products}.

Introduce a formal parameter \(z\). The componentwise
generating series of \eqref{eq:mc-derivatives} are
\[
 \sum_{r\geq1}z^r(\psi_r)_0(E)=D_E(x)-D_E(x-z),\qquad
 \sum_{r\geq1}z^r(\psi_r)_1(sf)=f(x-z)-f(x).
\]
For polynomial inputs these sums are finite. For formal
inputs they are identities of formal series in \(x,z\),
proved by the same binomial expansion. These formulas
identify the translation carried by the higher components.
They assert no numerical convergence at a prescribed
nonzero value of \(z\).
The translation acts on the coordinate of the potential.
Relating it to a worldsheet operation requires a geometric
comparison with its own coordinate data.

\section{Formal coefficients and the remaining comparisons}
\label{sec:mc-formal}

For formal power series, use the complete modules and
total-degree filtrations of Section~\ref{sec:da-formal}.
If the bulk operators have coefficient orders at least
\(A_1,\ldots,A_r\), and the boundary maps have orders at least
\(F_1,\ldots,F_m\), the mixed component has order at least
\[
 \max\left(0,\sum_{j=1}^r A_j+\sum_{i=1}^m F_i-(m+r-1)\right).
\]
There are \(m+r-1\) homotopies in each summand, each lowering
order by at most one. Multiplication adds the input orders,
while \(i,\pi\) preserve them. Each fixed pair of arities
therefore gives a jointly continuous map.
The geometric series in \eqref{eq:mc-lifted-maps} uses
completion in boundary arity. Its terms at a fixed arity
are finite compositions on the coefficient modules.

The map from diagonal operators to coherent central
operations now preserves their associative products through
all higher components. A general equivalence still requires
its full kernel and cokernel to be determined with the
appropriate critical-value support. Compatibility with
Hochschild brace operations, cyclic pairings and chiral
operations imposes additional equations beyond this
associative comparison.
